\section{Conclusiones}

El principal aprendizaje fue que I/O no bloqueante no significa reemplazar una
llamada bloqueante por otra con una bandera. Cada evento habilita sólo una parte
del trabajo: una lectura puede traer medio mensaje, una escritura puede aceptar
algunos bytes y un \texttt{connect} puede quedar pendiente. Modelar ese progreso
en estados y buffers hizo posible conservar datos y aplicar contrapresión sin
detener a las demás conexiones.

La resolución de nombres mostró otro aspecto menos evidente: sacar
\texttt{getaddrinfo} del hilo principal resuelve el bloqueo, pero introduce un
problema de identidad y tiempo de vida. Tuvimos que razonar qué ocurre si el
cliente se desconecta, el descriptor se reutiliza o el proceso comienza a
apagarse mientras la resolución continúa. Aprendimos que la concurrencia debe
analizarse por propiedad de los datos y no sólo por la cantidad de hilos.

Diseñar PMC también cambió nuestra forma de describir un protocolo. Los ejemplos
sirven para ilustrar, pero no reemplazan una gramática, una máquina de estados ni
reglas que delimiten cada respuesta. El prefijo de cantidad en respuestas
multilínea y el orden en pipelining surgieron de preguntarnos cómo podría un
cliente independiente interpretar el flujo sin conocer nuestra implementación.

Por último, las mediciones de carga nos enseñaron a separar una observación de
una afirmación de capacidad. Las gráficas permiten comparar throughput, latencia
y errores bajo las combinaciones ejecutadas; no demuestran por sí solas un máximo
universal. El límite de \texttt{pselect}, dado por \texttt{FD\_SETSIZE}, y el consumo de dos
descriptores por conexión explican por qué el entorno y la metodología forman
parte del resultado.
